\section{An Efficiency-Gated Verifiable Reward}
\label{sec:reward}

Section~\ref{sec:degeneracy} says a correctness-only reward cannot work. It
does not say what should replace it, and the space of replacements is not
empty - several are already published. This section states the one the
degeneracy argument implies, explains why its \emph{structure} rather than its
ingredients is the claim, and describes the suite we built to test it against
the alternatives.

\subsection{The gate}

For a program $p$ with correctness indicator $c(p)$ and reference speedup
$S_p(n)$ measured on $n$ threads, define the capped strong-scaling efficiency
\begin{equation}
  E_p(n) = \min\!\left(\frac{S_p(n)}{n},\ 1\right)
  \label{eq:efficiency}
\end{equation}
and the efficiency-gated reward
\begin{equation}
  \Rgate(p) = c(p)\cdot E_p(n).
  \label{eq:rgate}
\end{equation}

Equation~\eqref{eq:efficiency} is not new and we do not present it as new: it
is the strong-scaling efficiency the HPC community has used for decades, and
in the LLM setting it is essentially \ParEval's own $\mathrm{efficiency}_n@k$
metric~\cite{pareval} with a cap at unity. Naming that lineage matters, because
the point of the paper is not the metric but where it belongs - in the reward,
not only in the evaluation table.

Under $\Rgate$, an incorrect program scores $0$; the serial projection of
Definition~\ref{def:projection} scores approximately $Q_p/n$, which for a
program of ordinary sequential quality is $1/n$ - $0.125$ at eight threads;
and a program achieving linear scaling scores $1$. Proposition~\ref{prop:degeneracy}
no longer holds: $\serialproj(p)$ is now separated from $p$ by
$E_p(n) - 1/n$, which is the entire dynamic range of the reward.

\subsection{Why a gate and not a weighted sum}

The published correctness-plus-performance rewards for code are, with
consistency, \emph{additive}. ACECode~\cite{acecode} awards $0.5$ for passing
tests and adds up to $0.5$ of an efficiency term. Kevin~\cite{kevin} awards a
correctness constant plus a speedup ratio. Real-machine HPC
post-training~\cite{realmachine} uses a staged penalty ladder terminating in a
throughput term. Each is a reasonable engineering choice and each was validated
in its own setting.

The degeneracy argument implies a different structure, for two reasons.

First, an additive reward has a \emph{floor} that correctness alone reaches. A
program that is correct and entirely serial banks the correctness constant
unconditionally. The optimum is no longer serial, but the serial program is no
longer near-zero either; it is a safe, defensible answer that a
risk-averse policy can retreat to whenever parallelization is hard. A
multiplicative gate makes the performance term unreachable without correctness
\emph{and} makes correctness worth nothing without performance. There is no
floor to retreat to.

Second, and more practically, the gate preserves the property that made
verifiable rewards attractive in the first place: chasing reward can never buy
a wrong answer. Under an additive reward with a large enough efficiency
coefficient, a fast incorrect program can outscore a slow correct one, and
whether it does is a matter of hyperparameter tuning. Under
Equation~\eqref{eq:rgate} it cannot, ever, at any coefficient.

We state this as a structural argument, not an empirical superiority claim.
Section~\ref{sec:results-bakeoff} measures the difference; we did not want the
comparison to rest on the argument alone.

\subsection{What the gate does not fix}
\label{sec:reward-limits}

Adding a measured quantity to a reward adds a measurement to attack. We name
the attacks we know of, and we build them (Section~\ref{sec:results-hacks})
rather than only discussing them.

\paragraph*{Baseline manipulation}
If efficiency is computed from self-speedup $\Sigma_p(n) = T_p(1)/T_p(n)$, the
policy can raise its reward by making the one-thread run slower. This is the
same family as CUDA-L1's timing exploits, where added asynchronous streams
misled a main-stream timer~\cite{cudal1}. Anchoring to the fixed reference,
$S_p(n) = T_{R_t}/T_p(n)$, removes this: the denominator is not under the
policy's control.

\paragraph*{Algorithmic speedup in parallel clothing}
Anchoring to the reference introduces the dual problem. By
Equation~\eqref{eq:decomposition}, $S_p(n) = Q_p\Sigma_p(n)$, so a program that
improves the sequential algorithm by a factor of $n$ and uses no threads at all
reaches $E_p(n)=1$. The cap conceals it. We therefore report $Q_p$ and
$\Sigma_p(n)$ separately throughout, and recommend that any deployed reward do
the same; a reward that reports only $S_p(n)$ cannot distinguish a good
parallel program from a good sequential one.

\paragraph*{Timing-region evasion and caching}
Work moved outside the timed region, lazily evaluated, or memoized across the
repeat runs used for noise reduction will all inflate the efficiency term. The
mitigations are ordinary and unglamorous - time the whole driver, randomize
inputs per repeat, run repeats in fresh processes - and we specify them in
Section~\ref{sec:setup-protocol}.

\paragraph*{The measurement itself}
Unlike $c(p)$, $E_p(n)$ is a physical measurement with a noise distribution. A
reward whose between-program differences are smaller than its within-program
variance is not a reward, it is noise with a mean. Our protocol
(Section~\ref{sec:setup-protocol}) times in fresh processes and reports the
minimum over repeats to suppress this; quantifying a per-machine noise floor and
the resulting reward resolution is a measurement we recommend any deployed
performance reward report.

\subsection{\PG: an adversarial suite for parallel-code rewards}
\label{sec:reward-paragate}

A reward for parallel code should be judged the way a verifier is judged: by
what it accepts that it should not. We therefore built \PG, a suite of
\NHacks\ hand-written candidates spanning \ParEval's OpenMP problem families,
each of which is \emph{correct} and each of which defeats at least one natural
reward. The classes are listed in Table~\ref{tab:hackclasses}.

\begin{table}[t]
\caption{The \PG\ hack classes. Every candidate is functionally correct and
therefore earns the maximum correctness-only reward. The final column names
the reward property each is designed to defeat.}
\label{tab:hackclasses}
\centering
\footnotesize
\begin{tabular}{@{}llp{2.7cm}@{}}
\toprule
ID & Construction & Defeats \\
\midrule
H1 & every \texttt{\#pragma omp} deleted & correctness-only \\
H2 & \texttt{parallel for num\_threads(1)} & static ``has a pragma'' checks \\
H3 & loop body wrapped in \texttt{critical} & nominal parallelism \\
H4 & one-thread path artificially slowed & self-speedup $\Sigma$ \\
H5 & better serial algorithm, no threads & reference speedup $S$, capped \\
H6 & work hoisted outside the timed region & timer-scoped rewards \\
H7 & results cached across repeat runs & repeat-based noise reduction \\
H8 & correct but $O(n)$ oversynchronized & additive rewards with a floor \\
\bottomrule
\end{tabular}
\end{table}

The suite is the reusable part of this paper. A reward formulation can be
scored against it in minutes, and the resulting matrix
(Figure~\ref{fig:hackheat}, Section~\ref{sec:results-hacks}) is a far more
informative summary of a reward than the single benchmark number usually
reported. We release it with the intention that subsequent reward designs for
HPC code report against it, including designs that beat ours.

A suite, though, is only as sharp as the programs and timings it is scored over.
The next section fixes those: the tasks, the models that actually wrote the code,
the frozen pool every reward is computed on, and the timing protocol that makes
the efficiency term reproducible enough to serve as a reward rather than an
anecdote.
